\section{Introduction}
\label{sec:introduction}

\citet{shin2026compositional} introduced compositional self-improvement, a way
for a model to extend itself to problems it cannot yet solve without any new
labels. The model splits a hard problem into easier ones it already solves,
combines their answers into an answer for the hard problem, and trains on the
result. Repeating this takes it far beyond the problems it started with. The
appeal is plain: labelled data is scarcest exactly where a model is weakest,
and here the model supplies its own.

The same paper shows that each such training round can be a gamble. A round
might help, do nothing, or leave the model worse than before, and
\citet{shin2026compositional} explain why. The answers to the easier problems
may be wrong. Combining them may introduce errors even when they are right.
And the answers that can be built may never include what the model still needs
to learn. From these they draw three lessons: teach the easy skills well, find
where combining breaks, and check what the built answers can reach.

The three lessons identified by \citet{shin2026compositional} are very useful,
but difficult to put into practice for a model that is improving itself
autonomously. This is because each lesson relies on something the model does
not necessarily have access to without human intervention
(Table~\ref{tab:what-each-needs}). To find out why a round failed, the authors
compare the answers the model built with the correct ones. To keep wrong
answers out of training, they use a filter written for the task. When adding
two long numbers, for example, joining the sums of their halves goes wrong
whenever a digit is carried from one half into the other, so their filter
discards those problems. A model teaching \emph{itself} has no correct answers
to the hard problems, and no one to write an effective filter. In this paper,
we therefore ask whether it is possible to predict the outcome of a round
\emph{before it is run, using only the model's own answers}. We introduce
three checks that make this possible(Section~\ref{sec:teacher-and-student}), so
that a round that is not expected to help can be skipped. We also show that
the checks require less compute than the obvious alternative, training briefly
and peeking at early progress.

\begin{table}[t]
\centering
\caption{What each way of assessing a training round requires, and what it
  tells us. The first two rows are from \citet{shin2026compositional}; the
  third is proposed here. A rule means a statement, written by hand, of when
  combining gives a wrong answer.}
\label{tab:what-each-needs}
\small
\begin{tabular}{lccl}
\toprule
 & Needs correct & Needs a rule & \\
Tool & answers & for the task & What it tells us about the round \\
\midrule
Diagnosis & Yes & No & Why a round failed \\
Filter & No & Yes & Which built answers to discard \\
Checks introduced in this paper & No & No & What a round will do, before it is run \\
\bottomrule
\end{tabular}
\end{table}

Our approach is to see one model in two roles. As the \emph{teacher}, it
applies the split-and-combine procedure; as the \emph{student}, it answers the
same problem directly. Training makes the student imitate the teacher, so a
round can help only if the teacher's answers differ from the student's, are
correct, and include something the student lacks. The last two conditions
restate the failure modes identified by \citet{shin2026compositional}; the
first is new. For each of these three conditions, we introduce a check that
can be performed before the round is run, using only the model's own outputs
(Section~\ref{sec:teacher-and-student}). 

We test all three checks in the expertimental setting used by \citet{shin2026compositional}, with
every threshold fixed before the final evaluation (Section~\ref{sec:setup}).
Before each round our checks are used to predict whether the round will help,
do nothing or do harm. This prediction matched what the round went on to do in
nearly all cases, far more often than any fixed guess would, and flagged every
harmful round, for only a small fraction of the training round's cost.
Starting the round and peeking at early progress was no more accurate, cost
several times as much, and needs correct answers
(Section~\ref{sec:do-the-checks-predict}). A fourth check that we had also
specified in advance, on how well the model solves the easier problems, failed
to predict harm. 

Our experiments yield three further findings about training on self-built answers: 
on how to select them, on what the teacher can build, and on what the student can learn from them (Section 5).

The evidence presented in this paper comes from one small model and two
synthetic tasks, chosen because every outcome can be verified and harmful
rounds can be produced on purpose. Although the thresholds we use are specific
to each task, we expect the three checks, and the framing of self-improvement
as a teacher--student loop, to be useful in other settings.
Section~\ref{sec:beyond-synthetic-tasks} examines how this framing accounts
for self-improvement loops reported by others.

Finally, Section~\ref{sec:discussion-and-limits} discusses the implications of
our work for self-improvement loops, and the work planned to test the proposed
checks in more complex settings. If the checks carry over, a recursively
self-improving model need not run a round to learn its outcome: from its own
answers alone, it could tell when another round will help, when it will do
harm, and when no further round can help.